% !TeX root = ../../../main.tex
\subsection{Open-Source Models for Avoiding Potential Data Leakage}
\label{ch07:apd:open_source}

We evaluated potential data leakage arising from large language models (LLMs) pretrained on web-scale corpora that may include PubMed Open Access (PMOA) case reports.
Our safeguards target two avenues of leakage: (i) \emph{content overlap} with PMOA text, and (ii) \emph{task leakage} from using representations too close to raw prose.
Below we summarize the controls embedded in our methods and the additional empirical evidence gathered for this concern.

\subsubsection{Task Design Minimizes Leakage}

\textbf{Temporal annotations—not raw text—drive all evaluations.}
We transform narratives into chronological lists of \emph{(event, time)} tuples and operate on these derived annotations (future-masked where applicable).
Models are never scored on lexical similarity to the raw prose.

\textbf{Temporal reordering changes the learning problem.}
Events are reordered by extracted timestamps, breaking the original textual sequence seen during pretraining.
This creates a distinct modeling challenge: correct temporal ordering and causal dependencies must be inferred from structured tuples rather than reproduced from prose.
As shown in the main text, training/evaluating on time-ordered versus text-ordered sequences yields measurably different performance.

\textbf{Temporal metrics reduce any advantage from memorization.}
Evaluation emphasizes temporal reasoning—forecasting F1 at fixed horizons, pairwise concordance (c-index), and AULTC—rather than surface-form similarity.
Memorizing text order can even \emph{hurt} time-alignment metrics because clinical narratives often interleave retrospective commentary where textual order $\neq$ chronological order.

\subsubsection{Open-Source, PubMed-Free Baselines}

\textbf{Decoder baseline.} 
\emph{RedPajama-INCITE-7B-Instruct} is pretrained on English CommonCrawl (67\%), C4 (15\%), GitHub (4.5\%), Wikipedia (4.5\%), Books (4.5\%), arXiv (2.5\%), and StackExchange (2\%), explicitly \emph{excluding} PubMed. In our experiments it performs comparably to decoders that may have seen PMOA during pretraining.

\textbf{Encoder baselines.} 
We include \emph{BERT} (BooksCorpus + Wikipedia), \emph{RoBERTa} (CommonCrawl News, OpenWebText, Stories, Wikipedia), and \emph{DeBERTa} (\textit{RoBERTa}-style corpora).
None of these encoders include PubMed in their pretraining data.


\subsubsection{Post-Cutoff External Validation (Additional Evidence)}
To probe encoder leakage further, we evaluated \textit{BERT}, \textit{RoBERTa}, and \textit{DeBERTa} on 115 case reports published \emph{after} their pretraining cutoffs (post-2020).
Results mirror submission trends without degradation, which is inconsistent with a memorization hypothesis.


% !TeX root = ../../../main.tex
\begin{table}[htbp]
\caption{Post-cutoff evaluation on 115 case reports (forecasting F1 at 0h/24h/168h).}\ChSevenFloatLabel{ch07:tab:postcutoff_f1}
\ChSevenTableBodyStart
\begin{ChSevenFloatBody}

\centering
\footnotesize
\setlength{\tabcolsep}{6pt}
\renewcommand{\arraystretch}{1.1}
\begin{tabular}{lccc}
\hline
\textbf{Model} & \textbf{0h} & \textbf{24h} & \textbf{168h} \\
\hline
BERT     & 0.63 & 0.62 & 0.51 \\
RoBERTa  & 0.54 & 0.61 & 0.65 \\
DeBERTa  & 0.58 & 0.56 & 0.53 \\
\hline
\end{tabular}



\end{ChSevenFloatBody}
\end{table}


We also report competitive performance from \emph{OLMO-32B-Instruct} (fully documented training data) and parity from the PubMed-excluded \emph{RedPajama-INCITE-7B-Instruct} relative to decoders potentially exposed to PMOA, further indicating that prior PubMed exposure does not meaningfully advantage our temporally structured tasks.

The \textbf{key takeaways} are 
\begin{itemize}
    \item Using derived, time-ordered \emph{(event, time)} tuples and temporal metrics directly targets temporal reasoning and reduces susceptibility to content memorization.
    \item PubMed-free baselines and post-cutoff tests provide external checks
    \item Empirically, models without PubMed exposure achieve comparable performance, arguing against causal leakage into our forecasting, concordance, and AULTC evaluations.
\end{itemize}


%%%%%%%%%%%%%%%%%%%%%%%%%%%%%%%%%%%%%%%%%%%%%%%%
%%%%%%%%%Appendix J:  
%%%%%%%%%%%%%%%%%%%%%%%%%%%%%%%%%%%%%%%%%%%%%%%
% \clearpage
